\makeatletter
\def\input@path{{src/}} % Usamos preámbulo para cargar path con estilos
\makeatother

% \usepackage{src/yatt}
\documentclass{yatt}

\title{\nombredocu}
\author{\myFullName}

% Idioma principal: Español
\newif\ifspanish
\spanishtrue
%\spanishfalse

\begin{document}
\pagestyle{fancy}
\fancyhf{}
\renewcommand{\headrulewidth}{0pt}
\fancyhead{}
\fancyfoot{}
\fancyfoot[C]{\thepage}
% -------------------------
% PORTADA
% -------------------------
\frontmatter
\begin{titlepage}
    \begin{center}
        \includesvg[width=2cm]{assets/logos/LogoUCLM.svg}\par\vspace{1cm}
        \vspace{4cm}
        {\huge\bfseries \nombredocu\par}
        {\large \descripciondocu}
        \vspace{3cm}
        \par
        %------------------------------------------------
        %	Author name and information
        %------------------------------------------------
        \normalsize
        \shadowbox{
        \begin{minipage}{0.8\textwidth}
        \raggedright % Right align the text
            \ \\
            \textbf{Autor:} \myname \\\vspace{.15cm}
            \textbf{Fecha:} \today \\%\vspace{-.1cm}
            \ifCajaFirma
              \hrulefill
              \vspace{.25cm}\\
              \textbf{Firma:} \\
            \fi
            \vspace{.25cm} % Espacio para la firma
        \end{minipage}
      }
    \end{center}
    \vspace{2cm}
\end{titlepage}


% -------------------------
% TABLA DE CONTENIDOS (TOC)
% -------------------------
\tableofcontents

% -------------------------
% CONTENIDO PRINCIPAL
% -------------------------
\mainmatter
\chapter{Hola!}

\lipsum[1-4]


% -------------------------
% SECCIÓN BIBLIOGRÁFICA
% -------------------------

\newpage
\phantomsection  % Ojo necesario con hyperref.
\addcontentsline{toc}{chapter}{\bibname} % Añade la bibliografía al Índice de contenidos
\bibliographystyle{apacite}
\nocite{*} % COMENTAR SI NO QUIERES MOSTRAR REFERENCIAS NO CITADAS PERO INCLUIDAS EN BIB
\bibliography{bibliography} % Nombre del fichero .bib (sin extensión)

% -------------------------
% ANEXOS
% -------------------------
% \backmatter
\chapter{Anexos}

\lipsum[4-7]


\end{document}
